\cotefile{10}
\scanpage{89}

% letterhead
INSTITUT\\
DES HAUTES ÉTUDES SCIENTIFIQUES\\
BURES-SUR-YVETTE (S.-\&-O.)\\
928-48-53

\uline{Filtration} \uline{sur foncteur fibre} \uline{pour catégorie tensorielle}

Soit \deleted{$E$} $\mathcal M$ une catégorie tensorielle \deleted{sur} ``rigide \unsure{connexe}'' sur corps $k$, $S$ un schéma sur $k$, $F$ un foncteur fibre relatif à $S$

\[
\mathcal M\longrightarrow\text{faisceaux loc libres sur } S
\]

Une filtration de $F$ par des sous-foncteurs $F^{(p)}$ est dite \uline{admissible} si elle satisfait aux conditions

(0) Chaque $F^{(p)}(M)$ est loc. fact. direct dans $M$ ($\forall M\in\mathrm{Ob}\,\mathcal M$)

\[
\text{(i) Additivité :}\quad F^{(p)}(M\ast N)=F^{(p)}(M)\ast F^{(p)}(N)
\]

\[
\text{(ii) Multiplicativité}\quad F^{(p)}(M\otimes N)=\sum_{q+r=p}F^q(M)\otimes F^r(N)
\]

(iii) $F^p(\check M)$ est l'orthogonal de $F^{-p\unsure{+1}}\deleted{\unsure{(}}(M)$

(iv) Pour chaque $M$, la filtration $F^{(p)}(M)$ est loc. exh. et sép.

Considérons alors le foncteur

\[
M\leadsto\operatorname{gr}F(M)
\]

En vertu de (0), il est à valeurs dans les modules loc libres \inserted{(gradués)} sur $S$, (i) indique qu'il est additif, (ii) qu'il est multiplicatif, et (iii) qu'il commute à $\check{\phantom M}$.

\deleted{\unsure{B}} N.B. \deleted{Ses conditions impliquent} (iii) et (iv) impliquent $F^{(0)}(1)=\underline{1}$, $F^{(1)}(1)=0$, d'où $\operatorname{gr}F(1)=\underline{1}$ réduit au degré $0$.

Supposons $\mathcal M$ $1/2$ simple, alors $\operatorname{gr}F$ \uline{exact} (sinon il faut le mettre en condition \unsure{de|en|ds|dans} (i)).

Si on pose
% archive sheet number: 88
